% Original analytic teaching diagram, 2026-11-12.
% Chain of chapter 7.28: score -> Fisher information -> local KL quadratic form -> natural gradient,
% with the condition needed at each arrow (differentiation under the integral, finite second moment,
% integrable second-derivative control, positive definiteness). Exponential family branch: G = Hessian of A
% = Jacobian of the mean map; mean coordinates carry G^{-1}.
\definecolor{VB}{HTML}{7188A3}\definecolor{VBF}{HTML}{E5EAF0}
\definecolor{VG}{HTML}{829C8B}\definecolor{VGF}{HTML}{E8EEE9}
\definecolor{VY}{HTML}{B4A16F}\definecolor{VYF}{HTML}{F1ECDD}
\definecolor{VR}{HTML}{A7787D}\definecolor{VRF}{HTML}{F0E5E6}
\definecolor{VP}{HTML}{9588A5}\definecolor{VPF}{HTML}{ECE8F0}
\definecolor{VI}{HTML}{3B4252}\definecolor{VM}{HTML}{6B7280}\definecolor{VL}{HTML}{4C4D4F}
\begin{tikzpicture}[x=1mm,y=1mm,>=Stealth,font=\figurefont\fontsize{8}{10}\selectfont,
lab/.style={text=VM,font=\figurefont\fontsize{7.5}{9}\selectfont,align=center},
box/.style={draw=VB,fill=VBF,line width=.6pt,rounded corners=1.2mm,minimum height=8mm,minimum width=20mm,align=center,text=VI,inner sep=1.2mm},
gbox/.style={draw=VG,fill=VGF,line width=.6pt,rounded corners=1.2mm,minimum height=8mm,minimum width=20mm,align=center,text=VI,inner sep=1.2mm},
ybox/.style={draw=VY,fill=VYF,line width=.6pt,rounded corners=1.2mm,minimum height=8mm,minimum width=20mm,align=center,text=VI,inner sep=1.2mm},
arr/.style={-{Stealth},draw=VL,line width=.7pt}]
\node[box] (s) at (0,20) {得分\\$\bm{s}_{\bm{\theta}}=\nabla_{\bm{\theta}}\ln p_{\bm{\theta}}$};
\node[box] (g) at (40,20) {Fisher信息\\$\bm{G}=\mathbb{E}_{\bm{\theta}}[\bm{s}\bm{s}^{\top}]$};
\node[box] (k) at (80,20) {局部KL\\$\tfrac12\bm{\delta}^{\top}\bm{G}\bm{\delta}$};
\node[gbox] (n) at (80,-6) {自然梯度\\$-\bm{G}^{-1}\nabla_{\bm{\theta}}L$};
\node[ybox] (e) at (22,-6) {指数族\\$\bm{G}=\nabla^{2}A$};
\draw[arr] (s) -- node[lab,above=1mm]{二阶矩有限} (g);
\draw[arr] (g) -- node[lab,above=1mm]{二阶可积控制} (k);
\draw[arr] (k) -- node[lab,right=1mm]{$\bm{G}$正定\\固定半径} (n);
\draw[arr] (e) -- node[lab,right=1mm,pos=.45]{均值映射的\\Jacobian} (g);
\node[lab,anchor=north] at (0,13.5) {$\mathbb{E}_{\bm{\theta}}[\bm{s}]=\bm{0}$：\\积分与求导交换};
\node[lab,anchor=west] at (35,-6) {均值坐标中为$\bm{G}^{-1}$};
\end{tikzpicture}
